\section{Methodology} \label{Met}

In this section, we first present an overview of our analysis and then explain in detail the different steps of the methodology in the corresponding subsections.
In the current work we perform an exhaustive analysis of the hemispherical power asymmetry anomaly for intensity and polarization following the methodology of \cite{Planck:2019evm}, where the local-variance estimator is used to find dipolar-like features. The estimator is introduced in subsection \ref{subsection:3.3}, and essentially the output is the amplitude and direction of the local-variance dipole, which can be use as a tracer of the modulation.

To apply this analysis to the polarization data, we first need to construct an E-mode map. This is a non-local quantity, and ideally we need full-sky maps of the Q and U Stokes parameters to have an accurate reconstruction of the E map. However, since CMB observations are strongly contaminated by residual foregrounds in certain regions of the sky, we can not consider in the analysis those pixels excluded by the polarization confidence mask. To solve this problem, a useful approach is to carry out inpainting in the contaminated regions, a process that assigns new values to the excluded pixels such that they are more consistent with the rest of the CMB map. In particular, we have implemented a method based on \emph{Gaussian Constrained Realizations} (GCR) which is used to inpaint the regions masked by the polarization confidence mask in the Q and U data maps.
In subsection \ref{subsection:3.1} the inpainting method is described in more detail. 
We note that this technique requires the knowledge of the covariance matrix of the total signal present in the map. For the CMB signal, 
this can be obtained analytically given a cosmological model, whereas for the noise and systematics component, we compute it from the E2E noise simulations. However, to avoid overfitting (due to the lack of convergence of the estimated matrix because of the relatively small number of simulations), 
the simulations considered in the analysis must be independent from those used to compute the covariance matrix. Therefore we are forced to split the simulations in two halves, the first set used to estimate the covariance matrix and the second one to obtain the distribution of the amplitudes of the local-variance dipole and the corresponding p-value.\footnote{Note that the p-value is defined in an frequentist way as the fraction of simulations with an amplitude of the estimator equal or larger than the one obtained from the data.} Different partitions of the simulations are also considered in order to test the robustness of the results. In particular, subsection \ref{subsection:3.2} describes how the p-value is estimated taking into account the different data splits.

Another important step is to check the validity of the E-mode map constructed from the inpainted Q and U data maps. In particular, using the E2E simulations, we generate an extended customised mask by imposing a threshold in the pixel error of the reconstruction of the E-map. This mask is then used for the analysis of the polarization data (see subsection \ref{subsection:3.2} for details on how this mask is constructed).

Finally, section \ref{subsection:3.4} presents some tests to validate our methodology.
We check with the E2E simulations that the estimated amplitude and direction of the local variance dipole are not biased as a consequence of the inpainting process. We also give a comparison between the inpainting scenario and the one where E-modes are derived from a simple masking of the Q and U Stokes parameters. Furthermore, assuming a dipolar modulation model (eq.~\ref{DM}), we obtain some interesting results such as 
the sensitivity of the method taking into account the realistic Planck noise and systematics.

\subsection{Inpainting using Gaussian constrained realizations}\label{subsection:3.1}

One of the main issues when analyzing the CMB is how to deal with foreground emissions such as dust, free-free, or synchrotron. Although component separation algorithms are able to reduce these emissions significantly, some residuals are expected to remain near the Galactic plane and in the
locations of point sources, where the sky is strongly contaminated. 
The standard approach is to mask these regions, 
but, as mentioned before, this leads to other  difficulties.
This is especially the case when dealing with polarization data,
since masking Stokes parameters introduces an undesired mixing between E and B modes \cite{Tegmark:2001zv, Lewis:2003an}. This effect has an impact 
when reconstructing the E mode map, but it can be minimized by using an optimal inpainting technique. Inpainting consists on replacing the contaminated pixels by values which are somewhat consistent with the cleaned data, keeping as much as possible the statistical properties of the underlying field.
Different inpainting approaches have been considered in CMB analyses. For instance, 
diffuse inpainting, which consists on replacing the value of the contaminated pixels by averages of those of the neighbouring pixels, works acceptably well to inpaint small regions, such as those 
corresponding to point sources \cite{Planck:2018yye}. A second more sophisticated approach is purified inpainting, which constructs a full-sky polarization map from a masked one where most of the pure E- and B-modes are projected out in order to minimize the leakage. It has been used in the Isotropy and Gaussianity analyses performed by the Planck Collaboration \citep{Planck:2019evm}, including the study of the hemispherical power asymmetry.
In the present work, we use an alternative approach 
based on GCR. This technique was introduced in \cite{Marcos-Caballero:2019jqj} for temperature, but we have extended it here for the Stokes parameters.

GCR is a method that works in the pixel domain. The idea is to fill the masked pixels by sampling from the conditional probability distribution, $p(\mathbf{\hat{d}}|\mathbf{d})$, where $\hat{d}$ is the vector of the inpainted field and d is the vector of the available pixels. Assuming that the field is Gaussian, which is a very good approximation for the CMB, all we need is the pixel covariance matrix.
For the CMB signal, this 
can be computed from the theoretical power spectra\footnote{We use the $\Lambda$CDM best fit power spectra available in the PLA.} following, for instance, the Appendix A\footnote{
Note the small typo in equation A7 of the appendix
where the proportionality constant is negative, not positive. Thus, in equation A8, the negative sign corresponds to the case where z component of the vector $\hat{r}_{ij}$ is positive, and vice versa.} of \cite{Tegmark:2001zv}, whereas for the noise plus systematics part, we need to rely on simulations. We will present a full description of the method in a future work, together with a user-friendly
code (that will be publicly available), and some applications (Gimeno-Amo et al., in preparation). 
Although, in principle, this method is able to fill the contaminated pixels with a realisation that is perfectly consistent with the remaining data according to the assumed statistical properties of the underlying field, it presents two main limitations. First of all, this approach is very demanding from the computational point of view, since it requires the computation and storage of 
the pixel covariance matrix, whose dimension is $2N_{\rm pix} \times 2N_{\rm pix}$ when considering the inpainting of the Q and U Stokes parameters ($3N_{\rm pix} \times 3N_{\rm pix}$ if also TE correlation is considered). Hence the resolution at which we can work is
limited by the required memory. 
Due to this constraint, we have carried out our analyses at a Healpix resolution of $N_{\rm side}=64$. The second limitation comes from our lack of knowledge of the part of the covariance matrix coming from noise and systematics in the Planck polarization data, 
which needs to be characterised using simulations. In the case of the Sevem method for PR4, we have a total of 600 E2E simulations (600 for the CMB and 600 for the noise and systematics). As previously mentioned, to avoid overfitting\footnote{Overfitting arise from the non convergence of the covariance matrix. If we use the same simulations to estimate the matrix and perform the analysis, the matrix will have exactly the correlations given by the sample of simulations, and this will produce an overfitting when inpainting is applied on them.} 
in the inpainting process, we need to split them in two sets of 300 independent simulations, such that the first set is used to characterize the noise contribution to the covariance matrix, while the second one is used to obtain the distribution of the amplitude of the local variance dipole that will be compared with the data. 
The limited number of E2E noise simulations is not enough to achieve a good convergence of the noise covariance matrix, and consequently, matrix elements, especially the off-diagonal terms, are 
 significantly noisy. Indeed, using semi-realistic noise simulations -- Gaussian simulations that include anisotropic and correlated noise according to the noise covariance matrix obtained from the Planck noise E2E simulations -- we have found that the number of simulationes required to achieve convergence is at least of a few thousand. Another approximation is to ignore the --in general-- non-Gaussian behaviour of systematics, that is not taken into account in the impainting. However, in spite of these limitations, 
 we find that this proposed inpainting procedure is still very useful for our purpose.
Indeed, it is not necessary to reproduce the full statistical properties of the inpainted pixels in order to minimise significantly the E-to-B leakage due to incomplete sky.

\subsection{Performance of the inpainting, confidence mask, and estimation of the p-value}\label{subsection:3.2}

The performance of the inpainting procedure can be determined by comparing the exact E-mode map obtained from full-sky Q and U maps of the E2E CMB plus noise simulations, with that obtained from the same Q and U maps that have been inpainted in the region given by the Planck polarization confidence mask.
In particular, we compute the map of E-mode residuals $\Delta{E}$, as the difference between the exact and inpainted maps, for the pixels outside the mask for each of the inpainted simulations.
We also obtain the standard deviation of the residuals at each unmasked pixel ($\sigma_{\Delta E}$) as well as the one of the input E-mode map ($\sigma_E$). In the case of PR3, this procedure is repeated for the simulations propagated through the four component separation method, while for PR4 is done for Sevem. A relative error map $\delta{E}$ is then constructed as follows:
%
\begin{equation}
\delta{E}=\frac{\sqrt{\sum_m \sigma_{\Delta{E}}^2(m)}}{\sqrt{\sum_m \sigma_{E}^2(m)}}
\end{equation}
%
where $m$ runs over the different component separation methods used for PR3 or PR4. Note that for PR4, since only Sevem is considered, this is simply the ratio between the dispersion of the residuals and that of the input E map at each pixel.
Figure \ref{errormap} shows a comparison between the map of the relative error in the E-map that we get by applying inpainting to the PR4 data (left panel) and by directly using masked Q and U maps (right). 
It becomes apparent that the inpainting works quite well, especially for point sources and isolated regions outside the Galactic plane. Moreover, in the ideal case in which the covariance matrix is perfectly known, we have found that the performance is even better. 

\begin{figure}[h!]
    \centering
    \includegraphics[scale = 0.35]{Figuras/Projections_Error_1_new.png}
    \includegraphics[scale = 0.35]{Figuras/Projections_Error_2_new.png}
    \caption{Comparison of relative error maps (in percentage) in the reconstruction of the E-mode for PR4. \textit{Left panel: }Error map obtained with the inpainting approach. \textit{Right panel: }Error map obtained for the simple Q and U masking approach (no inpainting). Grey area corresponds to the PR3 polarization confidence mask.}
    \label{errormap}
\end{figure}

Despite having a more accurate reconstruction of the E-mode using inpainting, the residuals in regions close to the mask boundaries are still significant. One possibility to reduce these residuals is to extend the polarization confidence mask. In particular, we define a set of extended masks by selecting different thresholds in the error map, which are determined from a compromise between the maximum admissible error and the minimum 
extension of the mask. 
Figure \ref{FigMask} shows the variation of the maximum error with $f_{\rm sky}$\footnote{The sky fraction is actually determined from a combined mask obtained from the study of different splits. More specifically, the combined mask is obtained as the product of the individual extended masks generated for each of the considered splits (see below).} for PR4.
By removing the pixels more affected by the inpainting, an improvement in the maximum error is rapidly achieved, going down to values of $\delta{E}\sim$ 40\% with only a moderate reduction of the sky. However, further improvements require the removal of a significant fraction of the sky as a trade-off.
The complete set of masks will be used to check the robustness of the results against the considered sky fraction, but we select as reference mask the one with the threshold of $\delta{E}$ = 40\% because, for PR4, it leaves a fraction of sky similar to that of \cite{Planck:2019evm}.

This way of generating confidence masks is similar to the one followed in \cite{Planck:2019evm}, although with some differences. The first one is the resolution, since in that work the mask is generated at $N_{\rm side} = 1024$ and then downgraded to $N_{\rm side} = 64$, while we generate it directly at the final resolution (mainly due to our limitation to work with high resolution maps). 
The second difference comes from the way in which the threshold is chosen, since in their work the residuals of both the E and B-modes are considered.
We decided to focus only on the E-mode reconstruction for several reasons: B is consistent with zero in the Planck data, the E-mode signal is much larger than the B-mode, and finally we are only applying the estimator to E maps.

\begin{figure}[h!]
    \centering
    \includegraphics[scale = 0.35]{Figuras/Figura_Error_vs_Fsky.pdf}
    \caption{Sky fraction as a function of the allowed maximum error for the PR4 data set. 
    Vertical lines represent the thresholds for the different masks used in the analysis, covering from 45\% to 30\% error. The red star shows the maximum error and the fraction of the sky for our reference mask, which is used for the main results of this work and plotted in the right panel of Fig. \ref{Mask}.}
    \label{FigMask}
\end{figure}

Despite the increase in the error of the estimated dipolar modulation parameters due to the non-convergence of the noise covariance matrix, the amplitudes remain unbiased (see Section \ref{subsection:3.4}). However, splitting the simulations into two halves introduces additional uncertainty in the determination of the p-value. Then, the question is: how should we split the simulations? Although any random split is equally valid, each partition will yield a different p-value due to the slight variation in the amplitude distribution. There are also other effects that contribute to this dispersion but at a lower level. For example, the estimated noise covariance matrix is also different, which in turn affects slightly the way in which the inpainted is performed.

To get some additional insights on which is the best procedure to obtain the required p-value using the limited number of simulations, we have made some tests (see below) using semi-realistic noise simulations. These are constructed as Gaussian realisations characterised by the covariance matrix obtained from the full set of 600 E2E noise simulations\footnote{More specifically we calculate the Cholesky decomposition of the covariance matrix that is then used to generate the Gaussian realizations with the required statistical properties (see e.g. \cite{bar08}) for PR4.}. 
Therefore, these semi-realistic simulations of anisotropic noise contain the same correlations and statistical information up to second order as the realistic Planck simulations. The possible presence of non-Gaussianity is the only feature that we can not simulate. According to these simulations, a good estimation of the underlying p-value, the one that we would get if we sampled the distribution of amplitudes using the complete set of 600 simulations, can be achieved if we proceed as described below (this is detailed for PR4, but a similar procedure is applied to PR3):

\begin{enumerate}
    \item We generate 30 random splits from the 600 E2E noise simulations. For each split we have one pair of independent sets of 300 simulations each, i.e., a total of 60 sets in total.
    \item For each split, we use the first set
    to estimate the noise covariance matrix and we apply inpainting on the second set.
    We repeat the same procedure but this time estimating the matrix with the second set and applying it to the first one.
    \item For each of the estimated covariance matrices (a total of 60), we construct an inpainted map from the real data.
    \item For each of the 60 inpainted sets we generate an extended confidence mask by thresholding the residuals as explained previously for the reference case (i.e., maximum $\delta_E$ of 40 per cent). These are then combined into a single extended mask as the product of all of them. We repeat this procedure for different thresholds to have a set of extended confidence masks. Figure \ref{Mask} shows the reference extended masks for PR3 and PR4, used to obtain the main results of our analyses. 
    \item For each of the inpainted sets and inpainted data map (in each case inpainted using the same covariance matrix), and using the combined reference confidence mask, we apply the local-variance estimator (see section \ref{subsection:3.3}) to estimate the amplitude of the dipole. Finally, we define the p-value as the fraction of simulations with an amplitude equal or greater than the one observed in the data map.
    \item We get the mean of the 60 calculated p-values as our estimation of the underlying p-value. For the direction, we get the mean of 
    each of the vector components, and then, convert it to Galactic longitude and latitude coordinates.
\end{enumerate}

To establish the validity of the previous approach, the test we performed was the following one. First, we generate 2000 CMB simulations modulated in Q and U according to equation \ref{DM} using an amplitude and direction consistent with the ones estimated from the data. To each modulated CMB simulation, we add a semi-realistic noise simulation. Then we also generate 600 unmodulated CMB simulations to which we add 600 semi-realistic noise simulations (independent of the previous 2000).
Using this full set of 600 simulations, we estimated the p-value for each of the 2000 modulated simulations. This was carried out by applying the estimator to the ideal E-mode map, reconstructed from the full sky Q and U maps (i.e., without inpainting)\footnote{Although the E-maps are full-sky, we still apply an extended mask calculated following the same procedure as for the E2E simulations.}. We label these p-values as the \emph{true} ones, since they are the best possible estimation obtained from 600 simulations and without the uncertainties introduced by the inpainting. 
Once we have our reference p-values, we proceed as described above, i.e., we calculated the p-value from 30 different partitions of the 600 unmodulated simulations (including the inpainting process), obtaining a total of 60 p-values for each of the 2000 modulated simulations.
We find that the mean value of these distributions of 60 values was a good representation of the previously estimated \emph{true} value. As expected, the potential bias largery depended on the original p-value. In particular, the bias for the range of p-values between 0\% and 8\% was below 1/300. Although the p-values obtained from different splits are correlated, we provide, along with the mean value, the range of p-values containing 68\% of the distribution as a rough estimation of the uncertainty.

The same procedure has been followed for the PR3 data set, except that, in this case, we have only 300 E2E noise simulations to construct the 30 different partitions. Therefore for each split we have one pair of independent sets of 150 noise simulations each (used to construct the covariance matrix). In addition, we also have 999 CMB simulations, for which we consider 900 in our study. In this way, we construct sets of 450 simulations that are used for the analysis, obtained by combining one noise simulation with three different CMB simulations.

\begin{figure}[h!]
\centering
  \includegraphics[scale = 0.35]{Figuras/Customize_PR3_Mask.png}
  \includegraphics[scale = 0.35]{Figuras/Customize_PR4_Mask.png}
  \caption{\textit{Left panel: }Customized mask used for the PR3 data set analysis. It leaves 61.4\% of the sky unmasked. \textit{Right panel: }Customized mask used for the PR4 Sevem data set analysis. It leaves 64.6\% of the sky unmasked. The grey region corresponds to the polarization confidence mask provided by Planck, while black region corresponds to the extended area where the relative error in the E-mode reconstruction  
  is larger than the 40\%.}
  \label{Mask}
\end{figure}

\subsection{Local-Variance estimator}\label{subsection:3.3}

As previously mentioned, the estimator we have used to characterize the asymmetry is the Local-Variance introduced by Y. Akrami \cite{Akrami:2014eta}. The main motivation is that a dipolar modulation of the anisotropies would manifest itself as a dipolar structure in a map generated by computing variances over sky patches.
The procedure to obtain the amplitude and direction of the dipolar modulation is the following:
%
\begin{enumerate}
    \item This first step is exclusive to the analysis of polarization data due to the low signal-to-noise ratio. As mentioned in \cite{Planck:2019evm}, if we apply directly the method to the data, then the analysis returns local-variance dipoles whose directions are not uniformly distributed across the sky (see Figure 32 in \cite{Planck:2019evm}). Actually, their distribution is strongly correlated with the structure of the anisotropic noise. Therefore, we consider the following transformation directly applied on the input E-mode maps:
    %
    \begin{equation}
        \label{eq:Bias}
        X^{'}_{i} = (X_{i}-M_{i})/\sigma_{i}^{2} 
    \end{equation}
    %
    where $X_{i}$ is the value of the E mode map at pixel i computed from the full-sky inpainted Q and U maps. $M_{i}$ and $\sigma_{i}$ are the mean value and standard deviation at pixel $i$ obtained from inpainted simulations.

    \item We fix a specific HEALPix low resolution and we define a set of discs of a certain radius centered in the pixel centroids of the HEALPix map. Following previous works, we use $N_{side}$ = 16 and discs of 4 degrees for the present analysis. 
    
    \item We identify all the pixels of the map at the original resolution (in our case, $N_{side} = 64$ for polarization and $N_{side} = 64$ and 2048 for temperature)  that are inside each of the discs, compute their variance only taking into account the unmasked pixels, and associate these values to the pixel in the low resolution map. This is what we know as Local Variance Map (LVM). If more than 90\% of the pixels inside a disc are masked, then we mask the pixel in the LVM and remove it from the analysis. An example for one PR4 simulation is shown in Figure \ref{fig:LVM}.

    \item From the set of the LVM of the simulations, we estimate the mean $\bar{y}_i$ and the variance $\sigma_{d, i}^{2}$ on each disc, corresponding to one pixel at resolution $N_{\rm side}=16$. 

    \item Finally, we fit a dipole to each of the local-variance maps by applying a weighted $\chi^{2}$ 
    %
    \begin{equation}
        \label{chi}
        \chi^{2}=\sum_{i}\frac{[(y_{i}-\bar{y}_{i})-d_{0}-\textbf{d}\cdot \hat{\textbf{r}}_{i}]^{2}}{\sigma_{d, i}^{2}}
    \end{equation}
    %
    where $y_{i}$ is the LVM at pixel i; 
    $d_{0}$ captures the monopole component of the LVM; \textbf{d} = ($d_{x}$, $d_{y}$, $d_{z}$) is the dipole component; and $\hat{\textbf{r}}_{i}$ is the unit vector pointing to the i-th pixel. Sum is over all non-masked pixels in the LVM. $d_{0}$ and \textbf{d} are obtained by minimizing the $\chi^{2}$, which is analytic in this case.
\end{enumerate}


\begin{figure}[t!]
    \centering
    \includegraphics[width=\textwidth]{Figuras/LVM_Example.png}
    \caption{\textit{Left:} Sevem PR4 E2E E-mode simulation after applying the 
    transformation given by Eq.~\ref{eq:Bias}. The map is masked with an extended mask allowing a sky fraction of 63.2\%  \textit{Right:} Local-variance map generated following the previous steps. The grey region corresponds to those pixels discarded since more than 90 per cent of the pixels associated to the disc are masked. A disc of $4^{\circ}$ is used.}
    \label{fig:LVM}
\end{figure}

\subsection{Sensitivity and validation with simulations}\label{subsection:3.4}

Before going into the results, we assess the sensitivity of the method by 
considering CMB simulations that contain a dipolar modulation, taking also into account the realistic E2E PR4 noise and systematics simulations. 
In particular, we remark that for PR4, a transfer function that affects only the lowest multipoles of the polarization is present in the data (see section 4.3 and Fig. 20 from \cite{Planck:2020olo}). In order to take this effect into account for the E-map, we have computed the transfer function over the full-sky\footnote{Note that the transfer function is actually anisotropic, but using the full-sky transfer function is a sufficiently good approximation for our purpose.} for the Sevem data as: 
%
\begin{equation}
f^{EE}_{\ell} = \left( \frac{1}{N_{\rm sim}} \sum_{j=1}^{N_{\rm sim}} \frac{C_{\ell, j}^{\rm CMBx(CMB+N)}}{C_{\ell, j}^{\rm CMB}} \right)^2
\label{eq:TF}
\end{equation}
%
where $C_{\ell, j}^{\rm CMB}$ corresponds to the EE power spectrum of the $i$th PR4 Sevem CMB simulation and $C_{\ell, j}^{\rm CMBx(CMB+N)}$ to the cross-spectrum between that simulation and the same one after adding the associated noise and systematics Sevem simulation\footnote{The effect of the transfer function is actually included in the noise and systematics simulation, which presents therefore a correlation with its asssociated CMB simulation.}. Figure \ref{TF} shows the full-sky E-mode transfer function estimated for Sevem.

\begin{figure}[h!]
    \centering
    \includegraphics[scale=0.4]{Figuras/Transfer_Function_Sevem.pdf}
    \caption{Sevem PR4 full-sky E-mode transfer function obtained from equation \ref{eq:TF}.}
    \label{TF}
\end{figure}

Starting with the E2E Sevem CMB simulations, we construct modulated maps (modulated set) using Eq. (\ref{DM}). Each simulation is modulated for 14 different amplitudes (going from 5\% to 18\%) and using the same direction measured for the data (see section \ref{PR4_E_Results}). The modulation is applied to the Q and U maps, as a direct modulation in the E-mode map is not physically allowed \cite{Kothari:2018cmt}. The transfer function is then applied to the modulated CMB map\footnote{Starting at $\rm N_{side}$ = 2048, we downgrade the maps to $\rm N_{side}$ = 512 removing the Gaussian beam of 5$'$ and the pixel window function. Then, we modulate the Q and U maps, and we downgrade them to $\rm N_{side}$ = 64 applying to the $\rm a_{\ell \rm m}$ the square root of the Sevem E-mode transfer function and the Gaussian beam of 160$'$. We go through an intermediate resolution to avoid the numerical errors in the functions that convert from map to Fourier space.}. Finally, the E2E noise simulations are added to the simulated CMB signal\footnote{Note that the $i$th CMB simulation is added to an independent noise simulation (the $i$th+1 one) to avoid including the transfer function.}.
All the tests presented here have been done using the customized mask with $f_{\rm sky}$ = 63.2\%. Although the specific results may be slightly different using other masks, the validation and main conclusions of this sub-section still hold.

The results of this sub-section have been obtained by combining a number of splits of the data set. In particular, as previously explained, we consider 30 different splits, a total of 30 pairs of 300 E2E noise simulations. For a given split, we use one of the sets to compute the covariance matrix, and, after adding the corresponding CMB E2E simulation to the second set, we carry out the inpainting on it (unmodulated analysis set). Then, from the 600 modulated simulations we take those (half of them) whose associated noise simulations are not included in the set considered to estimate the covariance matrix. This is done to avoid again the overfitting problem. We apply inpainting on this modulated simulations using the same covariance matrix as the one used for the unmodulated analysis set, and for each simulation we compute the local-variance dipole amplitude and the corresponding p-value taking as reference the unmodulated analysis set. 
Therefore, we can repeat this process for each modulated simulation a total of 30 times\footnote{Note that for the case of the real data, since the data are independent of the simulations, we can obtain 2 p-values for each split (i.e. both data sets in the split can be ued to construct the covariance matrix) and thus obtain a total of 60 p-values.}. This means that, for each considered amplitude of the modulation, we have a total of 18000 p-values (30 splits x 600 simulations). Although these values are not independent, they provide us with additional information with respect to use a single split and allows us to construct an approximated distribution of the estimated amplitudes. 

The left panel of Figure \ref{Detectability} shows the distribution of the p-value of the estimated amplitude of the dipolar modulation obtained as previously explained. The results correspond to simulations modulated with an amplitude of 9\%, which is close to the value found for the data (see section \ref{PR4_E_Results}). According to this figure, we would be able to find a p-value lower than 1\% in 39.3\% of the cases, while this percentage increases to 61.7\% for a p-value lower than 5\%. Additionally, the right panel shows similar information but for all the considered amplitudes. In particular, the percentage of times that we get
a p-value equal to or lower than 1\%, 5\% and 32\% for each amplitude is given. For Planck PR4 we find a p-value lower than 1\% in at least 95\% of the times for amplitudes larger than 16\%. The amplitude reduces to 14\% if we consider a p-value lower than 5\%. 
Using the set of modulated simulations we can also calibrate the method to give a relation between the local variance and the dipolar modulation amplitudes (see section \ref{PR4_E_Results}).

\begin{figure}[t!]
    \centering
        \includegraphics[scale = 0.37]{Figuras/P_values_distribution_A_9_and_detectability.pdf}
    \caption{\textit{Left panel:} p-value distribution of the estimated amplitude of the dipolar modulation for the modulated simulations. We use an amplitude of 9\% and a fixed direction of ($\ell$, b) = ($235^{\circ}$, $-17.5^{\circ}$). \textit{Right panel:} Probability curves as a function of the dipolar modulation amplitude with the realistic PR4 noise and systematics. 
    Assuming a dipolar modulation in the Stokes Q and U parameters, we could claim a detection with a p-value lower than 1\% (dashed line) in 95\% of the cases if the amplitude was larger than 16\%. Solid (dotted) lines correspond to the probabilities of getting a p-value equal or lower than 5\% (32\%).}
    \label{Detectability}
\end{figure}

We have also checked how well the method performs in estimating the direction of the dipolar modulation
by using the modulated simulations with an amplitude of 9\% (selected taking into account the results from Section \ref{PR4_E_Results}). In this case, we consider the ideal case of Q and U full-sky maps, to study the performance of the method itself, without including possible distortions introduced by the inpainting technique. We compare the longitude and latitude measured in the local variance maps for the 600 modulated simulations with the input values (which are fixed to ($\ell$, b) = ($235^{\circ}$, $-17.5^{\circ}$)), and compute the angular distance between them. Figure \ref{AngularDistance_Sensitivity} shows the distribution of the angular distance, where bins are not uniformly distributed, but instead they are selected so the area covered by each of them is constant. As seen, the dispersion in the estimated direction is quite large, even in this ideal case, although one should take into account that the considered amplitude is relatively small.

\begin{figure}[h!]
    \centering
    \includegraphics[scale = 0.4]{Figuras/AngDist_Reconstruction.pdf}
    \caption{Distribution of angular distance between the input direction ($\ell$, b) = ($235^{\circ}$, $-17.5^{\circ}$) and the one obtained after applying the LVM estimator on the modulated (A = 9\%) PR4 simulations.  The ideal case, in which the Q and U full sky maps are known, is considered.}
    \label{AngularDistance_Sensitivity}
\end{figure}

We have carried out a final test to study the specific effect of the inpainting in both the estimated amplitude and direction of the dipolar modulation. For this, we have estimated those quantities from the 600 (unmodulated) E2E simulations in the ideal case in which Q and U full-sky are considered (that we call \emph{input} case). Then we have estimated the same quantities using inpainting. For comparison, a simple masking approach, where the masked area is simply replaced by zeroes in the Q and U maps, has also been considered. The left panel of Figure \ref{Bias_A_Dir_And_Dist} displays the distribution of the difference $\Delta{A}$ between the input amplitude and the one recovered from the inpainting or the simple masking approaches. It becomes apparent that the inpainting technique improves significantly with respect to simply masking the data. Indeed, on average, the latter tends to overestimate the amplitude, while no significant bias is found for the former. Furthermore, the dispersion of the distribution is significantly larger for masking (0.11) than the one we get when using our inpainting technique (0.04).

On the right panel, we also display the angular distances for the same scenarios. Again, we are not using uniform bins. From the figure it is clear that the inpainting provides a better reconstruction of the direction than the simple masking. We also note that although the dispersion introduced by the inpainting in the estimation of this quanitity is significant, it is clearly below the intrinsic dispersion of the method (see Fig.~\ref{AngularDistance_Sensitivity}). Furthermore, the median of the blue (inpainting) distribution is 10 degrees, while for the green one (masking) is almost 34 degrees. 

\begin{figure}[h!]
    \centering
    \includegraphics[scale = 0.42]{Figuras/Bias_Amp_and_AngDist_All_Splits.pdf}
    \caption{\textit{Left panel: }Distribution of $\Delta{A_{LV}}$, defined as the difference between the amplitude measured on the inpainted (blue) or masked (green) simulations (denoted as output) and the one obtained when full-sky simulations are used (denoted as input). 
    Red and black lines correspond to the mean of the blue (inpainting) and green (masking) distributions, respectively. \textit{Right panel: }Angular distances (in degrees) between the direction measured on the inpainted simulation and the input ones (blue). Input values are obtained from the corresponding full-sky (non-inpainted) simulations, so we remove the dispersion introduced by the estimator itself by making the difference. In green the same quantity is plotted for the case where a simple mask is applied on the Stokes parameters.}
    \label{Bias_A_Dir_And_Dist}
\end{figure}